% Appendix C — soft-mode escape methodology (H3As Im-3m -> R3m).
\section{Appendix C --- The soft-mode-escape methodology}
\label{app:C}

This appendix documents the campaign's most transferable contribution: a
procedure for converting a dynamically \emph{unstable} high-symmetry
candidate into a \emph{stable} lower-symmetry polymorph that retains strong
coupling. The worked example is \ce{H3As}.

\paragraph{The instability.}
\ce{H3As} in the high-symmetry $Im\bar{3}m$ structure is a dynamical saddle
point: a $4^3$-$q$ \texttt{ph.x} census finds imaginary modes ($25/192$
imaginary, minimum $-\SI{1267}{\per\centi\meter}$). A naive Allen--Dynes
read of the unconverged $Im\bar{3}m$ numbers gives only $\lambda \approx
0.97$, $T_c \approx \SI{30}{\kelvin}$ --- the lattice is collapsing before
strong coupling sets in.

\paragraph{The escape (method).}
\begin{enumerate}\itemsep0.25em
\item Identify the softest (most negative) phonon eigenvector at the
  unstable $q$-point. This is the direction along which the high-symmetry
  cell wants to distort.
\item Displace the structure along that eigenvector and relax. The symmetry
  drops --- here $Im\bar{3}m \to R3m$ (rhombohedral), a trigonal distortion
  of the cubic cell.
\item Re-run the full electron--phonon calculation in the distorted cell. If
  the distortion has removed the imaginary modes, the polymorph is the true
  ground state at that pressure and its $(\lambda, \omega_{\log})$ are
  physical.
\end{enumerate}

\paragraph{The result.}
The $R3m$ \ce{H3As} polymorph is dynamically stable ($14/14$ $q$-points
clear) and carries genuine strong coupling: $\lambda = 1.65$, $\omega_{\log}
= \SI{450}{\kelvin}$. The Allen--Dynes critical temperature is verified to
libm precision (g5):

\begin{lstlisting}[basicstyle=\ttfamily\footnotesize,frame=single,xleftmargin=0pt]
$ hexa verify --expr allen_dynes_tc 1.65 450 0.10 55.88 --tol 0.01
  calc   = 55.8806  ~= expected 55.88  (within tol 0.01, |D|=0.000567718)
  tier   = GREEN SUPPORTED-NUMERICAL  (H3As R3m soft-mode-escape polymorph)
\end{lstlisting}

\noindent The escape recovers a stable, strongly-coupled phase
($T_c = \SI{55.9}{\kelvin}$) from a candidate that the high-symmetry screen
would have discarded as ``unstable, weak.'' The coupling nearly doubled
($0.97 \to 1.65$) once the lattice was allowed to find its real ground state.

\paragraph{Why this matters beyond \ce{H3As}.}
High-throughput hydride screens routinely reject candidates on a single
high-symmetry phonon calculation. Soft-mode escape says that a negative
eigenvalue is an \emph{instruction}, not a verdict: it points to the
lower-symmetry polymorph that should be evaluated instead. The procedure is
generic --- it requires only the soft eigenvector, which \texttt{ph.x}
already produces --- and it is the methodological complement to the N5 wall
(Appendix~\ref{app:D}): the wall says strong harmonic coupling is often a
soft-mode artifact, and escape is the disciplined way to find out whether a
stable polymorph survives underneath.

\paragraph{Diagnosing the soft eigenvector: \ce{H3O} as a case study.}
The method's first step --- identifying which atoms move along the soft mode
--- can often be settled before any displacement, from symmetry and
mass-weighting alone. \ce{H3O} in the high-symmetry $Pm\bar{3}m$
inverse-perovskite (O at the $1a$ inversion centre, three H at $3d$
edge-midpoints) carries 43 imaginary branches across $16$ irreducible
$q$-points, the worst at the zone boundary $L$ ($-\SI{1133}{\per\centi\meter}$)
and a triply-degenerate $T_{1u}$ optical mode at $\Gamma$
($-\SI{682}{\per\centi\meter}$). Two independent arguments fix the
eigenvector on the hydrogen sublattice:
\begin{itemize}\itemsep0.2em
\item \emph{Mass-weighting.} Since $\omega\propto1/\sqrt{m}$ and
  $\sqrt{m_O/m_H}\approx4$, oxygen-dominated branches live below
  $\sim\SI{500}{\per\centi\meter}$; every imaginary branch here is in the
  H-dominated high-frequency band.
\item \emph{Group theory (decisive).} O sits at the inversion centre, so its
  displacements are absorbed entirely into the acoustic $T_{1u}$ translation;
  the $\Gamma$ optical $T_{1u}$ at $-\SI{682}{\per\centi\meter}$ is a
  \emph{pure} H-sublattice transverse displacement (proton off-centering /
  O--H bend), with O stationary.
\end{itemize}
The instruction is therefore clear: freeze the zone-boundary distortion (the
strongest channel) into a periodicity-doubled, lower-symmetry supercell and
relax --- exactly the soft-mode-escape procedure, with the caveat that the
relaxed low-symmetry phase typically lowers $\lambda$. \ce{H3O}'s harmonic
$\lambda_{\mathrm{BZ}}=2.5$, $T_c\approx\SI{190}{\kelvin}$ are extrapolated
\emph{through} the imaginary branches and are physically untrustworthy until
a stable polymorph is fixed; the anharmonic SSCHA band (9--\SI{109}{\kelvin})
is the honest envelope (Appendices~\ref{app:B}, \ref{app:D}).

\paragraph{Three escape routes (generality).}
Soft-mode escape is one of three ways an imaginary mode can be cured, and the
\ce{H3O} diagnosis shows when each applies:
\begin{enumerate}\itemsep0.2em
\item \textbf{Lower-symmetry relax} (the \ce{H3As} route): follow the soft
  eigenvector to the true ground-state polymorph. Definitive when a stable
  polymorph exists underneath.
\item \textbf{Pressure} (proton-well restoration): compression can turn a
  double-well proton potential back into a single well, hardening the soft
  mode --- the mechanism behind the \ce{H3Br} pressure scan
  (Appendix~\ref{app:F}).
\item \textbf{Anharmonic SSCHA}: a harmonic imaginary curvature can become
  real once anharmonic free-energy curvature is included (the \ce{H3S},
  \ce{LaH10} precedent). Note an isotope swap alone does \emph{not} cure a
  harmonic imaginary mode --- the curvature is mass-independent.
\end{enumerate}

\paragraph{Honest scope.}
\ce{H3As}-$R3m$ at \SI{56}{\kelvin} is not a record and not ambient. The
methodology, not the number, is the contribution. The $R3m$ verdict is a
campaign-internal DFT result pending an independent atlas-registration
record; the $T_c$ value itself is \tierGreen\ closed-form, but the structural
identification carries the usual DFT-level (not measured) status (R4).
